\documentclass[11pt]{article}

\usepackage[margin=1in]{geometry}
\usepackage[T1]{fontenc}
\usepackage[utf8]{inputenc}
\usepackage{lmodern}
\usepackage{amsmath,amssymb,mathtools}
\usepackage{booktabs,longtable,tabularx,array}
\usepackage{enumitem}
\usepackage[numbers,sort&compress]{natbib}
\usepackage[colorlinks=true,linkcolor=blue,citecolor=blue,urlcolor=blue]{hyperref}

\setlength{\parskip}{0.6em}
\setlength{\parindent}{0pt}
\setlist[itemize]{leftmargin=*}
\setlist[enumerate]{leftmargin=*}

\newcolumntype{L}[1]{>{\raggedright\arraybackslash}p{#1}}
\newcommand{\code}[1]{\texttt{#1}}

\title{Inputs and Outputs of Major Codes in a Stellarator Design Workflow}
\author{stellarator\_workflow repository draft}
\date{March 11, 2026}

\begin{document}
\maketitle

\begin{abstract}
This document isolates the software contracts that appear in a modern
stellarator design loop. For each major code family it lists the main input
artifacts, the main output artifacts, the subset of outputs that are normally
handed to the next code in the pipeline, and the subset of outputs that are
commonly used as optimization or design objectives. The emphasis is on the
practical interfaces that connect VMEC or DESC equilibrium calculations to
Boozer transforms, neoclassical solvers, gyrokinetic solvers, and global
transport tools \citep{HirshmanWhitson1983,DudtKolemen2020,Boozer1983,Landreman2014SFINCS,Mandell2024GX,BanonNavarro2023}.
\end{abstract}

\tableofcontents

\section{Purpose and conventions}

The goal here is narrower than the main workflow manuscript. This document is
about \emph{interfaces}: what each code expects to read, what it writes or
returns, which outputs are consumed by the next stage, and which outputs are
typically minimized, constrained, or monitored during stellarator optimization.

Three conventions are used throughout:

\begin{itemize}
  \item ``Input'' means either an on-disk file accepted by the public CLI or a
  public in-memory object accepted by the documented Python or JAX API.
  \item ``Output'' means either a persisted artifact such as \code{wout\_*.nc},
  \code{boozmn\_*.nc}, \code{sfincsOutput.h5}, or a documented result object
  returned by a public API.
  \item ``Optimization objective'' means a quantity that is commonly included
  directly in an objective function or used as a screening metric, even when it
  is \emph{not} passed to the next code in the transport loop.
\end{itemize}

The repository snapshots under \code{\_sources/} were used to verify the
representative NetCDF and HDF5 schemas quoted below, so the tables are tied to
actual artifacts rather than only to README summaries
\citep{VMECRepo,VMECJaxRepo,DESCRepo,BoozXformRepo,BoozXformJaxRepo,NEORepo,NEOJaxRepo,SFINCSRepo,SFINCSJaxRepo,MONKESRepo,GXRepo,SpectraxGKRepo,TrinityDocs,NEOPAXRepo}.

\section{Executive summary}

\subsection{Pipeline hand-offs}

{\small
\begin{longtable}{L{0.14\linewidth}L{0.19\linewidth}L{0.28\linewidth}L{0.29\linewidth}}
\caption{What each stage usually hands to the next stage.}
\label{tab:handoffs}\\
\toprule
Code family & Main output artifact & Next code(s) & Subset usually handed forward \\
\midrule
\endfirsthead
\toprule
Code family & Main output artifact & Next code(s) & Subset usually handed forward \\
\midrule
\endhead
VMEC / vmec\_jax &
\code{wout\_*.nc} or in-memory wout-like object &
BOOZ\_XFORM, booz\_xform\_jax, GX, Trinity3D, NEOPAX &
equilibrium spectra and profiles, especially \code{rmnc}, \code{zmns},
\code{lmns}, \code{bmnc}, \code{bsubumnc}, \code{bsubvmnc}, \code{iotas},
\code{phi}, \code{phipf}, geometry scalars, and pressure / current profiles \\

DESC &
\code{*.h5} equilibrium family, plus optional exported VMEC or Boozer files &
MONKES directly; otherwise the same downstream codes as VMEC after export &
surface and equilibrium coefficients, pressure / iota / current profile
parameters, and any exported \code{wout}- or \code{boozmn}-compatible file \\

BOOZ\_XFORM /\\ booz\_xform\_jax &
\code{boozmn\_*.nc} or in-memory Boozer object &
NEO, NEO\_JAX, SFINCS, sfincs\_jax, MONKES, NEOPAX &
Boozer field spectrum and radial profiles, especially \code{bmnc\_b},
\code{ixm\_b}, \code{ixn\_b}, \code{iota\_b}, \code{buco\_b},
\code{bvco\_b}, \code{phi\_b}, \code{phip\_b}, and the selected surface list \\

NEO / NEO\_JAX &
\code{neo\_out.*}, \code{neolog.*}, optional current / diagnostic files, or
in-memory surface results &
usually no direct transport consumer &
usually none; optional current-related outputs may be fed into outer current or
equilibrium updates \\

SFINCS / sfincs\_jax &
\code{sfincsOutput.h5}, optional \code{.npy} state or transport-matrix dumps &
Trinity3D or an outer equilibrium / transport loop &
most often \code{particleFlux\_vm\_rN} and \code{heatFlux\_vm\_rN}; if
\code{includePhi1} is enabled in the Trinity3D adapter, the hand-off becomes
\code{particleFlux\_vd\_rN} and \code{heatFlux\_vd\_rN}; bootstrap-current and
\code{E\_r} information may also be reused \\

MONKES &
in-memory \code{Dij}, \code{f}, \code{s}, or an HDF5 monoenergetic database &
NEOPAX &
for today's NEOPAX reader the effective hand-off is the reduced database:
\code{D11}, \code{D13}, \code{D33}, \code{Er}, \code{Er\_tilde},
\code{drds}, \code{rho}, and \code{nu\_v} \\

GENE / GENE-3D &
installation-specific output files or restart / diagnostic datasets &
Trinity3D-like profile solvers, validation workflows &
usually species heat and particle fluxes; in linear studies only
\(\gamma\) and \(\omega\) may be kept \\

GX &
\code{run\_name.out.nc}, \code{run\_name.nc}, \code{run\_name.big.nc},
\code{run\_name.restart.nc} &
Trinity3D &
steady estimates of particle and heat flux, read in the present Trinity3D
adapter from either \code{Diagnostics/ParticleFlux\_st} and
\code{Diagnostics/HeatFlux\_st} or from \code{Fluxes/pflux} and
\code{Fluxes/qflux}; zeta-resolved stellarator runs can also use
\code{ParticleFlux\_zst} and \code{HeatFlux\_zst} \\

SPECTRAX-GK &
in-memory diagnostics and optional CSV / plot outputs &
outer optimization or a transport wrapper &
the natural hand-off is the same turbulent contract as GX:
\code{gamma}, \code{omega}, heat-flux traces, and particle-flux traces \\

Trinity3D &
\code{*.nc}, \code{*.bp}, \code{*.log.npy}, \code{*.out} &
VMEC or DESC in the next outer iteration &
updated density and temperature profiles, hence an updated pressure profile;
in some workflows bootstrap-related current information is also recycled \\

NEOPAX &
in-memory \code{Lij}, \code{Gamma}, \code{Q}, \code{Upar}; solved profile
states; optional HDF5 outputs &
VMEC or DESC in the next outer iteration &
updated \code{E\_r(r)}, neoclassical fluxes, and pressure or current-related
profiles from the reduced transport solve \\
\bottomrule
\end{longtable}
}

\subsection{Outputs commonly used as objectives}

{\small
\begin{longtable}{L{0.16\linewidth}L{0.42\linewidth}L{0.23\linewidth}}
\caption{Outputs commonly used as stellarator optimization or design objectives.}
\label{tab:objectives}\\
\toprule
Code family & Outputs commonly used in objectives & Also a direct transport input? \\
\midrule
\endfirsthead
\toprule
Code family & Outputs commonly used in objectives & Also a direct transport input? \\
\midrule
\endhead
VMEC / vmec\_jax &
aspect ratio, volume, \(\beta\), rotational-transform targets, magnetic-well /
Mercier-related metrics, boundary smoothness, equilibrium existence / residual
quality &
indirectly; the equilibrium itself is the upstream state for later codes \\

DESC &
quasisymmetry error, geometry constraints, aspect ratio, volume, rotational
transform, force-balance residuals, differentiable proxy objectives &
indirectly; exported equilibria feed the rest of the workflow \\

BOOZ\_XFORM / booz\_xform\_jax &
non-target Boozer harmonics, symmetry-breaking spectrum, mirror terms,
quasisymmetry error in Boozer coordinates &
usually no; mainly a geometry diagnostic and interface layer \\

NEO / NEO\_JAX &
\(\epsilon_{\mathrm{eff}}\), class-resolved effective ripple, trapped-particle
screening metrics, current proxies &
usually no \\

SFINCS / sfincs\_jax &
neoclassical particle flux, heat flux, bootstrap current, ambipolar
\(E_r\), \(\Phi_1\), transport matrices &
yes \\

MONKES &
monoenergetic \(D_{ij}\) coefficients and their sensitivity to radius,
collisionality, and \(E_r\) &
yes, through reduced neoclassical models such as NEOPAX \\

GENE / GENE-3D &
linear growth rates and frequencies, nonlinear heat and particle fluxes,
transport stiffness &
nonlinear fluxes yes; linear rates usually no \\

GX &
linear \(\gamma\) and \(\omega\), nonlinear heat flux, particle flux, free
energy, stiffness &
yes for heat and particle flux \\

SPECTRAX-GK &
\code{gamma}, \code{omega}, \code{heat\_flux}, \code{particle\_flux},
\code{Wg}, \code{Wphi} &
yes once wrapped into a transport loop \\

Trinity3D &
\code{Pfus\_MW}, \code{Qfus}, \code{Paux\_MW}, transport-consistent
\(n(r)\), \(T(r)\), and \(\beta\) &
it is already the global transport stage \\

NEOPAX &
ambipolar \(E_r\), bootstrap current, toroidal current, neoclassical fluxes,
reduced-order profile predictions &
yes \\
\bottomrule
\end{longtable}
}

\section{Code-by-code reference}

\subsection{VMEC and vmec\_jax}

\paragraph{Role.}
VMEC solves the fixed-boundary ideal-MHD equilibrium; \code{vmec\_jax}
preserves essentially the same physical contract while exposing a JAX-native,
differentiable API \citep{HirshmanWhitson1983,VMECRepo,VMECJaxRepo}.

\paragraph{Primary inputs.}
The canonical VMEC input is \code{input.NAME}. The core design variables are
the boundary Fourier coefficients \code{RBC(m,n)} and \code{ZBS(m,n)}, plus
pressure coefficients \code{AM} or the auxiliary pressure arrays
\code{AM\_AUX\_*}. Either rotational transform or current is prescribed through
\code{AI}, \code{AC}, and their auxiliary arrays. The magnetic scale is set by
\code{PHIEDGE}. Resolution and solver controls include \code{NS},
\code{MPOL}, \code{NTOR}, iteration budgets, and tolerance arrays. The JAX
version accepts the same input file and additionally supports in-memory resume
state and solver options.

\paragraph{Primary outputs.}
The main output is \code{wout\_*.nc}. Verified example files contain scalar
geometry and convergence quantities such as \code{aspect}, \code{Aminor\_p},
\code{Rmajor\_p}, \code{volume\_p}, \code{betatotal}, \code{b0},
\code{volavgB}, \code{fsqr}, \code{fsqz}, and \code{fsql}. The radial-profile
outputs include \code{presf}, \code{pres}, \code{phi}, \code{phipf},
\code{chi}, \code{chipf}, \code{iotas}, \code{iotaf}, \code{q\_factor},
\code{jcuru}, \code{jcurv}, \code{buco}, and \code{bvco}. The spectral
geometry is stored in arrays such as \code{rmnc}, \code{zmns}, \code{lmns},
\code{bmnc}, \code{gmnc}, \code{bsubumnc}, \code{bsubvmnc},
\code{bsubsmns}, \code{currumnc}, and \code{currvmnc}. The JAX API also
returns a \code{FixedBoundaryRun} object containing the solver state, flux
profiles, and a wout-like in-memory representation.

\paragraph{Subset handed to the next code.}
BOOZ\_XFORM and \code{booz\_xform\_jax} need the full equilibrium spectrum and
profiles encoded in \code{wout\_*.nc}. GX, Trinity3D, and NEOPAX geometry
readers also consume \code{wout}-level data rather than only high-level
scalars, because they need field-line geometry, the rotational transform, and
surface metrics.

\paragraph{Subset used as objectives.}
Common direct objectives are aspect ratio, volume, \(\beta\), target
\(\iota(s)\), and equilibrium-quality metrics. The residual scalars
\code{fsqr}, \code{fsqz}, and \code{fsql} are mainly QA signals rather than
physics design objectives.

\subsection{DESC}

\paragraph{Role.}
DESC is a differentiable pseudo-spectral equilibrium and optimization suite
\citep{DudtKolemen2020,Panici2023,Conlin2023,Dudt2023QS,DESCRepo}.

\paragraph{Primary inputs.}
DESC accepts an input file or Python configuration specifying the surface
representation, field period, toroidal flux, pressure profile, and either iota
or current profile, together with spectral resolution, continuation settings,
and optimizer options. It can also import VMEC-like inputs and outputs.

\paragraph{Primary outputs.}
The standard serialized output is \code{*.h5} or \code{*.pkl}. Example HDF5
equilibrium families contain groups such as \code{\_equilibria/*} with the
spectral coefficients \code{\_R\_lmn}, \code{\_Z\_lmn}, \code{\_L\_lmn},
surface and axis descriptions, toroidal flux \code{\_Psi}, and saved pressure
and iota profile parameters under \code{\_pressure/\_params} and
\code{\_iota/\_params}. DESC can also export VMEC-style or
BOOZ\_XFORM-compatible files through its VMEC utilities.

\paragraph{Subset handed to the next code.}
If the next code expects VMEC input, the hand-off is an exported \code{wout}.
If the next code expects Boozer data, DESC can provide a \code{boozmn}-style
file or a direct in-memory field object. MONKES can consume DESC-derived field
objects directly, while VMEC-centric tools usually rely on exported
\code{wout}-compatible data.

\paragraph{Subset used as objectives.}
DESC is itself commonly the optimization host, so quasisymmetry error, aspect
ratio, volume, iota targets, boundary regularity, and force-balance residuals
are often objectives. Once coupled to GX or another turbulent model, turbulent
transport metrics become part of the same objective stack.

\subsection{BOOZ\_XFORM and booz\_xform\_jax}

\paragraph{Role.}
These tools transform a VMEC-like equilibrium into Boozer coordinates
\citep{Boozer1983,BoozXformRepo,BoozXformJaxRepo}.

\paragraph{Primary inputs.}
The legacy route is a VMEC \code{wout\_*.nc} file plus an \code{in\_booz.*}
control file specifying resolution and the list of radial surfaces to compute.
The JAX version also accepts an in-memory wout-like object.

\paragraph{Primary outputs.}
The standard file output is \code{boozmn\_*.nc}. Verified examples contain
\code{iota\_b}, \code{pres\_b}, \code{beta\_b}, \code{phip\_b},
\code{phi\_b}, \code{bvco\_b}, \code{buco\_b}, \code{jlist}, mode lists
\code{ixm\_b} and \code{ixn\_b}, and Fourier arrays \code{bmnc\_b},
\code{rmnc\_b}, \code{zmns\_b}, \code{pmns\_b}, and \code{gmn\_b}. The JAX
API also holds Boozer \(I\) and \(G\) profiles and the selected radial grid in
memory.

\paragraph{Subset handed to the next code.}
NEO and NEO\_JAX need the Boozer spectrum and its radial profiles. SFINCS and
\code{sfincs\_jax} use the same Boozer geometry when running in a Boozer-based
geometry mode. MONKES and NEOPAX use the Boozer spectrum through direct field
or file readers.

\paragraph{Subset used as objectives.}
The most common design outputs are the non-target Boozer harmonics in
\code{bmnc\_b}, or equivalent symmetry-breaking measures built from them. These
are geometry objectives, not state variables advanced by a transport solver.

\subsection{NEO and NEO\_JAX}

\paragraph{Role.}
NEO computes effective ripple and related trapped-particle diagnostics from
Boozer data \citep{Nemov1999,Nemov2008,NEORepo,NEOJaxRepo}.

\paragraph{Primary inputs.}
The standard inputs are a \code{boozmn\_*.nc} file and a control file such as
\code{neo\_param.ext} or \code{neo\_in.ext}. The control file sets the surface
list, angular resolution (\code{theta\_n}, \code{phi\_n}), Fourier cutoffs,
Monte Carlo controls, accuracy targets, and switches for current calculation
and optional extra file dumps. NEO\_JAX also exposes a direct Python or JAX API
that returns structured result objects.

\paragraph{Primary outputs.}
The legacy outputs are \code{neo\_out.*} and \code{neolog.*}, plus optional
\code{neo\_cur.*}, \code{current.dat}, \code{conver.dat},
\code{diagnostic.dat}, \code{diagnostic\_add.dat}, and
\code{diagnostic\_bigint.dat}. The principal physics quantities are
\code{epstot} \((=\epsilon_{\mathrm{eff}}^{3/2})\), \code{epspar},
\code{reff}, \code{iota}, \code{b\_ref}, \code{r\_ref}, \code{ctrone},
\code{ctrtot}, \code{bareph}, \code{barept}, and \code{yps}. In NEO\_JAX
these appear in result objects under names such as
\code{epsilon\_effective} and \code{epsilon\_effective\_by\_class}.

\paragraph{Subset handed to the next code.}
Usually none. In most workflows NEO is a screening or optimization branch, not
a transport branch. If \code{CALC\_CUR=1}, current-related outputs can be fed
into a broader outer loop, but that is not the default Trinity3D hand-off.

\paragraph{Subset used as objectives.}
\(\epsilon_{\mathrm{eff}}\) and its class-resolved pieces are classic
stellarator screening objectives. This is precisely the case in which a code's
output is central to optimization but not central to profile evolution.

\subsection{SFINCS and sfincs\_jax}

\paragraph{Role.}
SFINCS solves the radially local drift-kinetic equation for neoclassical
transport in non-axisymmetric plasmas \citep{Landreman2014SFINCS,Mollen2018,SFINCSRepo,SFINCSJaxRepo}.

\paragraph{Primary inputs.}
The main input is \code{input.namelist}. In current public interfaces the key
groups are \code{general}, \code{geometryParameters},
\code{speciesParameters}, \code{physicsParameters},
\code{resolutionParameters}, \code{otherNumericalParameters},
\code{preconditionerOptions}, and optional \code{export\_f} settings. Together
they define the geometry source, radial coordinate, species charges and masses,
\(\hat n_s\), \(\hat T_s\), their gradients, the collision model,
electric-field terms, \(\Phi_1\) switches, and numerical resolution.

\paragraph{Primary outputs.}
The standard output is \code{sfincsOutput.h5}. Verified example files include
geometry snapshots, copied input metadata, and the following reusable transport
datasets:
\begin{itemize}
  \item \code{particleFlux\_vm\_psiHat}
  \item \code{particleFlux\_vm\_psiN}
  \item \code{particleFlux\_vm\_rHat}
  \item \code{particleFlux\_vm\_rN}
  \item \code{heatFlux\_vm\_psiHat}
  \item \code{heatFlux\_vm\_psiN}
  \item \code{heatFlux\_vm\_rHat}
  \item \code{heatFlux\_vm\_rN}
  \item additional transport state: matching momentum-flux arrays, classical
  fluxes, \code{FSABjHat}, \code{FSABFlow}, \code{Phi1Hat}, and
  \code{transportMatrix}.
\end{itemize}
Optional exports include full-\(f\), \(\delta f\), and \code{*\_vs\_x}
lineouts. The JAX interface additionally exposes in-memory result dictionaries
and can write \code{.npy} state vectors or transport matrices from its CLI.

\paragraph{Subset handed to the next code.}
The present Trinity3D adapter reads \code{particleFlux\_vm\_rN} and
\code{heatFlux\_vm\_rN} when \code{includePhi1} is off, and
\code{particleFlux\_vd\_rN} and \code{heatFlux\_vd\_rN} when it is on. In
other workflows \code{FSABjHat}, ambipolar electric-field information, and the
transport matrix are also reused by an outer iteration.

\paragraph{Subset used as objectives.}
Neoclassical particle flux, heat flux, bootstrap current, ambipolar \(E_r\),
and \(\Phi_1\)-related quantities are all common optimization metrics. Unlike
\(\epsilon_{\mathrm{eff}}\), many of these are also direct transport inputs.

\subsection{MONKES}

\paragraph{Role.}
MONKES is a fast JAX monoenergetic drift-kinetic solver
\citep{Escoto2024MONKES,MONKESRepo}.

\paragraph{Primary inputs.}
The public API constructs a field object from a DESC equilibrium or an
equivalent magnetic representation, then solves at a chosen radius with a
species list, radial electric field, particle speed, and pitch-angle
resolution. In database-building use, the same solve is repeated over
collisionality and \(E_r\) grids.

\paragraph{Primary outputs.}
The core output is the monoenergetic transport matrix \code{Dij}, together with
the perturbed distribution function \code{f} and source array \code{s}. In a
database context these are stored as HDF5 arrays such as \code{D11},
\code{D13}, \code{D31}, \code{D33}, \code{rho}, \code{nu\_v}, \code{Er},
\code{Er\_tilde}, \code{Es}, and geometric conversion factors.

\paragraph{Subset handed to the next code.}
NEOPAX's public database reader currently consumes the reduced subset
\code{D11}, \code{D13}, \code{D33}, \code{Er}, \code{Er\_tilde},
\code{drds}, \code{rho}, and \code{nu\_v}. The full distribution function
\code{f} and source term \code{s} are useful for analysis but are not part of
that hand-off.

\paragraph{Subset used as objectives.}
MONKES is often used to build cheap neoclassical surrogates. In that role the
\code{Dij} coefficients themselves are optimization or screening targets.

\subsection{GENE and GENE-3D}

\paragraph{Role.}
GENE and GENE-3D provide high-fidelity gyrokinetic turbulence calculations in
flux-tube and global settings \citep{Gorler2011GENE,Merz2012GENE,Xanthopoulos2020GENE3D,GENEWebsite}.

\paragraph{Primary inputs.}
The stable physics contract is clear even though the exact file syntax is not
publicly standardized across installations: geometry derived from VMEC or
Boozer coordinates, species densities and temperatures, radial gradients,
collision parameters, electromagnetic switches, and numerical-grid settings.
For stellarator flux-tube workflows the geometry usually passes through a
stellarator interface layer before reaching GENE.

\paragraph{Primary outputs.}
The usual outputs are linear growth rates, real frequencies, eigenfunctions,
and nonlinear species heat and particle fluxes together with spectra and time
histories. The exact filenames are installation-dependent, so the important
contract is the variable-level one rather than a specific disk layout.

\paragraph{Subset handed to the next code.}
For transport coupling, the critical hand-off is the turbulent flux vector. For
screening, only linear \(\gamma\) and \(\omega\) may be retained.

\paragraph{Subset used as objectives.}
Linear growth rates are often used for rapid pruning; nonlinear heat flux,
particle flux, and stiffness are higher-fidelity design objectives.

\subsection{GX}

\paragraph{Role.}
GX is a GPU-native gyrokinetic code with a workflow that is especially suitable
for repeated turbulent-transport evaluations inside design loops
\citep{Mandell2018,Mandell2024GX,GXRepo}.

\paragraph{Primary inputs.}
GX is run with \code{gx run\_name.in}. The input file controls geometry,
species, domain, time stepping, diagnostics, restart behavior, and numerical
resolution. Stellarator calculations can use VMEC-driven geometry or imported
field-line geometry, and the diagnostics flags \code{omega=true} and
\code{fluxes=true} determine whether growth-rate and flux diagnostics are
written.

\paragraph{Primary outputs.}
Linear runs write \code{run\_name.out.nc}. Nonlinear runs write
\code{run\_name.nc}. Saved-field diagnostics may also produce
\code{run\_name.big.nc}. Restart output uses \code{run\_name.restart.nc}.
Verified NetCDF outputs contain groups such as \code{Grids},
\code{Geometry}, \code{Diagnostics}, and \code{Inputs}. The
public docs show linear growth-rate diagnostics under
\code{Special/omega\_v\_time}. The present Trinity3D adapter reads flux data
from \code{Diagnostics/ParticleFlux\_st},
\code{Diagnostics/HeatFlux\_st}, their zeta-resolved variants, or from the
\code{Fluxes/pflux} and \code{Fluxes/qflux} groups when present.

\paragraph{Subset handed to the next code.}
Trinity3D uses steady estimates of the heat and particle fluxes at each radius.
It obtains flux Jacobians by rerunning GX on perturbed gradients and finite
differencing those outputs.

\paragraph{Subset used as objectives.}
Ion and electron heat flux are common direct objectives. Linear
\(\gamma\) and \(\omega\) are also used for screening. In turbulence-aware
optimization, heat flux is both an objective and a transport input.

\subsection{SPECTRAX-GK}

\paragraph{Role.}
SPECTRAX-GK is a JAX-native gyrokinetic solver aimed at differentiability and
rapid experimentation \citep{SpectraxGKRepo}.

\paragraph{Primary inputs.}
The current public interface is config-driven. A TOML config specifies grid,
geometry, physics toggles, and time integration. Geometry can be analytic or
imported from a GX-style \code{*.eik.nc} file, including VMEC-derived geometry.

\paragraph{Primary outputs.}
The natural outputs are JAX-native result objects with \code{gamma},
\code{omega}, and time traces such as \code{gamma\_t}, \code{omega\_t},
\code{Wg\_t}, \code{Wphi\_t}, \code{Wapar\_t}, \code{heat\_flux\_t}, and
\code{particle\_flux\_t}. The nonlinear CLI can also save a CSV with compact
columns for time, growth rate, frequency, free energy, and species-resolved
heat and particle flux.

\paragraph{Subset handed to the next code.}
The natural downstream contract is the same one used by GX: turbulent heat and
particle flux, and optionally linear \(\gamma\) and \(\omega\). In the current
public repository this is primarily an optimization or validation interface
rather than a fixed Trinity3D adapter.

\paragraph{Subset used as objectives.}
Growth rate, real frequency, turbulent heat flux, particle flux, and free
energy are all exposed directly and are therefore natural differentiable
objective components.

\subsection{Trinity3D}

\paragraph{Role.}
Trinity3D is a global transport solver that evolves density and temperature
profiles and computes whole-device power metrics \citep{TrinityDocs,BanonNavarro2023}.

\paragraph{Primary inputs.}
The main input is a TOML-style \code{*.in} file with groups such as
\code{[grid]}, \code{[time]}, one or more \code{[[model]]} blocks,
\code{[[species]]}, \code{[geometry]}, \code{[physics]}, and \code{[log]}.
Geometry can be read from a VMEC \code{wout}. Turbulence and neoclassical
models add their own inputs, for example \code{gx\_template} and
\code{gx\_outputs} for GX coupling.

\paragraph{Primary outputs.}
Trinity3D writes NetCDF, ADIOS2, NumPy-log, and text-log outputs. Verified
NetCDF outputs contain geometry groups with \code{B0}, \code{Btor},
\code{a\_minor}, \code{R\_major}, \code{area}, and \code{grho}; profile
histories such as \code{n\_e}, \code{T\_e}, \code{n\_H}, \code{T\_H}; total
and model-labelled fluxes \code{pflux\_*} and \code{qflux\_*}; gradients
\code{aLn\_*}, \code{aLT\_*}; source histories \code{Sn\_*} and \code{Sp\_*};
and whole-device metrics such as \code{beta\_vol}, \code{Paux\_MW},
\code{Palpha\_int\_MW}, \code{Pfus\_MW}, and \code{Qfus}.

\paragraph{Subset handed to the next code.}
The essential outer-loop hand-off is the updated profile state:
\(n_s(r)\), \(T_s(r)\), and therefore \(p(r)\). That pressure profile is what
must be sent back to VMEC or DESC if the loop is to be closed
self-consistently. When current evolution is modeled, bootstrap-related current
information can also be recycled into the next equilibrium solve.

\paragraph{Subset used as objectives.}
\code{Pfus\_MW}, \code{Qfus}, \code{Paux\_MW}, \(\beta\), and the final
transport-consistent profiles are end-of-pipeline design metrics.

\subsection{NEOPAX}

\paragraph{Role.}
NEOPAX is a JAX framework for reduced neoclassical transport and profile
evolution \citep{NEOPAXRepo}.

\paragraph{Primary inputs.}
Public examples build NEOPAX from three inputs: a VMEC \code{wout}, a Boozer
\code{boozmn}, and a MONKES-like monoenergetic database. The database reader
currently expects HDF5 datasets such as \code{D11}, \code{D13}, \code{D33},
\code{Er}, \code{Er\_tilde}, \code{drds}, \code{rho}, and \code{nu\_v}. The
profile state can be loaded from NTSS-like HDF5 files with arrays such as
\code{r}, \code{Er}, \code{Te}, \code{ne}, \code{Pressure}, \code{I\_bs}, and
related transport quantities, or assembled directly in memory.

\paragraph{Primary outputs.}
The core neoclassical API returns \code{Lij}, \code{Gamma}, \code{Q}, and
\code{Upar}. Example HDF5 outputs range from a minimal \code{Er\_Test.h5}
containing \code{rho}, \code{Er}, and \code{Jboots}, to larger NTSS-style
profile files containing arrays such as \code{Pressure}, \code{FluxNeo},
\code{FluxQe}, \code{FluxQi}, \code{AmbiFlux}, \code{J\_bs}, \code{I\_bs},
\code{I\_tor}, \code{beta}, \code{Te}, \code{TD}, \code{Tt}, \code{ne},
\code{nD}, and \code{nHe}.

\paragraph{Subset handed to the next code.}
In an outer loop the practical hand-off is the reduced transport state:
\code{Er(r)}, neoclassical fluxes, and any updated pressure or current
profiles. NEOPAX also contains turbulence-coupling utilities, but the public
examples center on the neoclassical reduced model.

\paragraph{Subset used as objectives.}
Ambipolar electric field, bootstrap current, toroidal current, neoclassical
fluxes, and reduced-order profile predictions are all candidate objective
components or screening metrics.

\section{The two most important distinctions}

Two distinctions repeatedly matter in stellarator workflow design:

\begin{enumerate}
  \item Some outputs are \emph{screening objectives only}. The clearest example
  is NEO's \(\epsilon_{\mathrm{eff}}\): it is central to ranking candidate
  geometries, but it is usually not what Trinity3D advances in time.
  \item Some outputs are \emph{both objectives and state hand-offs}. The
  clearest examples are turbulent heat and particle flux from GX, GENE, or
  SPECTRAX-GK, and neoclassical flux from SFINCS. These quantities are often
  minimized in optimization and are also the direct inputs needed to evolve
  density and temperature profiles.
\end{enumerate}

That distinction is what determines whether a future containerized black-box
workflow needs a result only for ranking, or whether it needs a robust,
machine-readable interface that another code must parse automatically.

\bibliographystyle{plainnat}
\bibliography{references}

\end{document}
